\chapter{Survey of the Literature}

This thesis approaches the problems of computation in the brain and in machine learning through understanding the rich spatiotemporal dynamics of visual searching. We use the stochastic FNS model, studying the properties of anomalous diffusion and complex multiscale spatiotemporal dynamics it yields, with good theoretical background in the neural circuit model it was developed for. We aim to  understand these behaviours as an efficient method of gathering information. In Sec~\ref{sec:stoch_modelling} we introduce the modelling framework of stochastic calculus and Lévy stable motion which informs the dynamics of the FNS sampler. Section~\ref{sec:comp_sampling} then provides the link between stochastic dynamics and computation sampling, both in biological and abstract contexts. Section~\ref{sec:behavioural_data} provides the background for much of the behavioural data we present and model in this thesis, and the previous methods used to understand and model the complex dynamics found in gaze trajectories especially.

\section{Stochastic Modelling}\label{sec:stoch_modelling}
Stochastic calculus has been fundamental in understanding how high dimensional complex systems reduce to a small number of variables of interest. When studying complex systems in physics, economics, geology, climate science, machine learning, and neuroscience, it is often the most natural, tractable, and accurate modelling choice when understanding a system of interest. The important concept here is that of the introduction of some central limit theorem, where the sum of very many small interactions which are understood microscopically can produce certain well-defined stochastic patterns on macroscopic scales of interest. The ubiquitous stochastic process is a Wiener Process $W_t$, equivalently called Brownian Motion. A Wiener Process is defined by the stationary and independent Gaussian increments across all non-overlapping time intervals. The Wiener Process is a very simple and analytically elegant model of the coarse-grained effect of very many small finite-variance interacting influences. A Wiener Process is defined only by its timescale, with a variance over trajectories after time $t$ of $t$ itself. Wiener Processes yield time series which are almost surely continuous but nowhere differentiable. The Wiener Process is understood as a model of the long-time dynamics of a variable effected by very many small uncorrelated finite variance influences \cite{hassler_stochastic_2016}. Classical examples of processes modelled by the Wiener Process are: the movement of a macroscopic particle in a non-turbulent fluid \cite{hassler_stochastic_2016}, the gravitational influence of an orbiting thermal collection of stars on supermassive black holes \cite{Chandrasekhar1943, Merritt2007a}, and (arguably) the movement of low-volatility stock prices \cite{cunchala2024basicoverviewvariousstochastic}. There are well studied mechanisms through which complex real-world stochastic processes diverge from the Wiener Process. Of interest to us is the case where the microscopic effects become correlated with heavy-tailed statistics. Correlated steps, long-time memory effects, and microscopic dynamics can result in a stochastic process which does not limit to the Wiener Process \cite{Brevitt2025, Klages2013}. Of interest to us is the fractional and superdiffusive Lévy flight, which is a result of spatial non-locality in an interacting network. For instance, previous studies have demonstrated a neural circuit network at criticality resulting in diverging correlation length of activity, and the centre of mass of neural firing is subject to the large jumps. Such large jumps in a stochastic time series are modelled according to Lévy stable motion, where the increments are drawn independently from a heavy-tailed distribution. This section describes the analytics and behaviour of such Lévy processes for our current study.
\subsection{Lévy Stable Motion}\label{subsec:levy_motion}
The literature is often unclear when distinguishing stochastic processes of a heavy-tailed nature. However, for our purposes, we analyse what is often called the \textit{Lévy flight} $L_T^{\alpha}$, a Stochastic Process defined through the increments being drawn from the heavy tailed symmetric $\alpha$ stable distribution (a $\mathcal{S}\alpha\mathcal{S}$ distribution). $\mathcal{S}\alpha\mathcal{S}$ distributions are a family of limiting stable distribution under summation of i.i.d infinite-variance variables, as a result of the generalised central limit theorem. The $\alpha$ value defining the particular $\mathcal{S}\alpha\mathcal{S}$ distribution defines the asymptotic tail of the probability density function, with $p(x) \propto |x|^{1-\alpha}$. Thus, the Lévy flight $L_t^{\alpha}$ after time $t$ is defined as follows, to outline the fractional scaling properties and alpha-stable statistics
\begin{equation}
    L_t \stackrel{d}{=} t^{1/\alpha} L_1, \ \ L_1 \sim \mathcal{S}\alpha\mathcal{S}
\end{equation}
where the parameter $\alpha$ defines the heaviness of the tails of resulting statistics. In the limit $\alpha \rightarrow 2$, the Lévy flight reduces to a scaling of Brownian Motion $L_t^{\alpha \rightarrow 2} = \frac{1}{\sqrt{2}}W_t$. If $1 < \alpha < 2$ the variance of $L_t^{\alpha}$ diverges, but $L_t^{\alpha}$ has 0 mean, whereas if $0 < \alpha \leq 1$, then both the variance and mean of $L_t^{\alpha}$ diverge. Many of the simple analytic properties of the Wiener Process are no longer true of a Lévy flight. Importantly, trajectories of a Lévy flight are almost surely no longer continuous. Instead, such trajectories are called \textit{càdlàg} (continue à droite, limite á gauche) with respect to time, which describes functions which are right continuous with well-defined left limits \cite{van_den_Bosch}. Thus, we say that since all such trajectories are càdlàg, Lévy flights are called a càdlàg process. Many useful properties of the càdlàg Lévy flight process are understood. It can be shown that Lévy flights are semimartingales, and may thus be integrated in the senses both of Ito and Stratonovich. Importantly, it follows from the theorem of the Lévy-Itô decomposition that a Lévy flight can be decomposed as the sum of two processes \cite{Kyprianou2014}:
\begin{enumerate}
    \item A compound Poisson process of large jumps with $|\Delta x| > 1$. Such jumps are exponentially distributed in time, and there are almost surely a finite number of such jumps in any finite time period.
    \item A square-integrable martingale jump process with $|\Delta x| < 1$. This process is thus finite-variance, and when coarse-grained over time, approaches a Wiener Process.
\end{enumerate}




\subsection{Anomalous Diffusion}\label{subsec:anom_diff}
\section{Abstract Computation}\label{sec:comp_sampling}
Probabilistic sampling is a powerful and flexible paradigm of general computation. Sampling methods represent and evaluate probability distributions through the production of stochastic samples samples. Uncertain, high-dimensional, and ambiguous inference problems require efficient and accurate characterisations of target distributions. Sampling methods draw from a target distribution such that the sampled ensemble distribution reflects the target itself. This paradigm has become widely utilised in machine learning and statistics, where generative models learn patterns from large multivariate distributions inferred from observational data \cite{delacour_neural_1997, cranmer_advances_2023}. Recently, this framework has gained significant traction as a model for biological computation. The position that the brain itself acts as a Bayesian generative model has become that interprets through continuous sampling mechanisms has become well supported.
\subsection{Active Searching and Inference}

\subsection{Neural Sampling}\label{subsec:neur_sampling}
The neural sampling hypothesis posits that the variability of a cortical neural response is an important functional feature of the biological computational sampling. Previously regarded as noise, the neural sampling framework interprets instantaneous neural activity patterns as stochastic samples drawn from a posterior distribution over possible sensory interpretations \cite{orban_neural_2016, NIPS2002_a486cd07}. As such, the location and frequency of firing patterns across cortical neurons are interpreted as stochastically sampling an informed target posterior, acting as the basis for some kind of computation. Consequently, perceptual uncertainty is directly encoded by the variability of cortical responses across time and across multiple trials, with higher uncertainty reflected in wider response distributions \cite{NIPS2002_a486cd07, delacour_neural_1997, orban_neural_2016}. Neural sampling thus provides an interpretation for the way the brain navigates uncertainty at the fundamental scale of computation. Various biological mechanisms have been proposed to explain how neural circuits generate samples. Traditional models often assume the existence of explicit stochastic noise sources, however recent models demonstrate that strongly recurrent, balanced inhibitory-excitatory deterministic neural networks can generate the high-dimensional variability which stochastic models marginalise over through chaotic dynamics. These emergent chaotic dynamics facilitate Bayesian learning, as for instance described by Terada and Toyoizumi through training a recurrent network to sample static features and dynamic trajectories, generalising learned priors to novel, uncertain environments through biologically plausible learning \cite{terada_chaotic_2024}. Furthermore, energy-based models (EBMs) have been developed which suggest that fast-spiking interneurons can estimate the local partition function to enable fast, localised sampling dynamics \cite{dong_neural_nodate}. The sampling hypothesis has thus given a convincing theory of biological neural computation when contrasted with the trial-averaged perspective of some averaged response combined with noise linked to neural variability \cite{orban_neural_2016}. However, traditional Gaussian or Poisson noise models representing these stochastic network dynamics inherently correspond to continuous Brownian motion and struggle to capture the full dynamics of a complex interconnected inhibitory-excitatory biological neuronal network \cite{koren_efficient_2025, kilgore_biologically-informed_2025, song_training_2016}. In reality, neural activity exhibits structure across multiple spatial and temporal scales and exhibit heavy-tailed, non-Gaussian statistics. This observation motivates the design and theory of the theory of Fractional Neural Sampling (FNS). FNS describes the complex spatiotemporal dynamics of spiking neural circuits, where localised activity patterns propagate not through ordinary diffusion, but via superdiffusive Lévy flights. Such dynamics are called fractional, as they are the result of nonlocal interactions requiring fractional calculus to describe \cite{tarasov_general_2021}. These which can be mathematically characterised by stochastic differential equations with fractional order derivatives. The critical advantage of FNS over standard continuous random walks is its ability to efficiently sample complex, multimodal probability distributions with far-apart modes—a major challenge when processing natural environments replete with multiple distinct salient features. While Brownian-motion-based samplers typically remain trapped in local modes due to high potential energy barriers, the heavy-tailed, large jumps inherent in FNS enable the neural sampler to adaptively and freely switch between multiple modes \cite{neal_sampling_1996, latuszynski_mcmc_2025, simsekli_fractional_nodate, metzler_random_2000}. This dynamical switching provides a mechanistic explanation for the bimodality of estimate distributions observed in human perceptual switching and visual perception inference.
\subsection{Markov Chain Monte Carlo Methods}\label{subsec:MCMC}
Markov Chain Monte Carlo (MCMC) methods constitute the algorithmic foundation for generating samples from target distributions in Bayesian inference and machine learning. MCMC operates by constructing a Markov chain whose equilibrium distribution is the desired posterior distribution. In the context of neural computation, biological circuits are often modelled as implementing continuous-time variants of MCMC, most notably Langevin sampling and Hamiltonian Monte Carlo (HMC). Langevin dynamics utilise local gradient information from the probability landscape, combined with Gaussian white noise, to explore the state space over time.  Despite their ubiquity, traditional MCMC methods encounter severe theoretical and practical limitations when dealing with multimodal distributions. Standard Langevin and HMC dynamics are driven by continuous Brownian and Hamiltonian motion respectively, and the traversal of low-probability regions, equivalent to dynamical energy barriers, becomes exponentially slow. Consequently, the mean exit time for the sampler to leave one mode and discover another grows exponentially with the modal separation distance. To accelerate this sluggish exploration, algorithms often introduce auxiliary momentum variables. However, even with momentum-based acceleration, traversing deep energy valleys remains highly inefficient \cite{latuszynski_mcmc_2025, neal_sampling_1996}. The fractional dynamics introduced in FNS-style sampling are one area of research into how fast mixing of samples in complex distributions can be achieved without the need for nontrivial parameter tuning. For instance, algorithms based on Hamiltonian Monte Carlo sampling and second-order Langevin sampling very similar to FNS dynamics have been briefly developed and explored \cite{simsekli_fractional_nodate, ijcai2018p419, qi_fractional_2022}. Under FNS dynamics, the mean exit time between modes scales linearly rather than exponentially with separation, demonstrating a profound computational advantage for efficient biological Bayesian inference \cite{qi_fractional_2022}.
\section{Behavioural Data}\label{sec:behavioural_data}
\subsection{Neural Frequency Bands}\label{subsec:neur_freq_bands}
Throughout much of this thesis, links are drawn between timescales of interactions and dynamics in the FNS model and well-studied frequency bands of activity in the human brain.
\subsection{Gaze Trajectories}\label{subsec:gaze_traj}
Healthy human gaze trajectories consist of different dynamical regimes. The most studied property of the pattern of gaze trajectories are their division into periods of fixation, interspersed by near-instantaneous jumps between fixations, these jumps being known as ``saccades'' \cite{erkelens_coordination_2006, seung_how_1996, lencastre_dynamical_2025, benitezy_eyetracking_1970}. Importantly, this behaviour of local movements interspersed with large discontinuous jumps suggest modelling eye movements through a Lévy flight. This method has a limited history in the literature, and fitting data to a simple phenomenological Lévy flight parametrised by the stable parameter $\alpha$ found limited use as a biomarker for ADHD \cite{papanikolaou_levy_2024}. Saccade metrics in general have been found to be valuable as non-invasive biomarkers for broader neural disorders, including schizophrenia, autism, ADHD, OCD, bipolar disorder, Parkinson's disease, and multiple sclerosis \cite{morita_eye_2017, bittencourt_saccadic_2013,zhu_diagnosis_2024, caldani_saccadic_2017, wang_saccadic_2025}. Periods of fixation have received less scholarly attention as a biomarker of cognition. Gaze fixations exhibit three distinct dynamical patterns: near-instant microsaccades within the same point of focus, slow ocular drift about a point of focus, and high frequency ocular micro tremor \cite{van_ede_looking_2026,graham_ocular_2023, engbert_integrated_2011, spauschus_origin_1999}. Examinations and literature reviews of these behaviours have often not considered their computational role in active sensing and efficient search.
